%-----------------------------------------------------------------------------------------------------------------------------------------------%
%	The MIT License (MIT)
%
%	Copyright (c) 2021 Jitin Nair
%
%	Permission is hereby granted, free of charge, to any person obtaining a copy
%	of this software and associated documentation files (the "Software"), to deal
%	in the Software without restriction, including without limitation the rights
%	to use, copy, modify, merge, publish, distribute, sublicense, and/or sell
%	copies of the Software, and to permit persons to whom the Software is
%	furnished to do so, subject to the following conditions:
%	
%	THE SOFTWARE IS PROVIDED "AS IS", WITHOUT WARRANTY OF ANY KIND, EXPRESS OR
%	IMPLIED, INCLUDING BUT NOT LIMITED TO THE WARRANTIES OF MERCHANTABILITY,
%	FITNESS FOR A PARTICULAR PURPOSE AND NONINFRINGEMENT. IN NO EVENT SHALL THE
%	AUTHORS OR COPYRIGHT HOLDERS BE LIABLE FOR ANY CLAIM, DAMAGES OR OTHER
%	LIABILITY, WHETHER IN AN ACTION OF CONTRACT, TORT OR OTHERWISE, ARISING FROM,
%	OUT OF OR IN CONNECTION WITH THE SOFTWARE OR THE USE OR OTHER DEALINGS IN
%	THE SOFTWARE.
%	
%
%-----------------------------------------------------------------------------------------------------------------------------------------------%

%----------------------------------------------------------------------------------------
%	DOCUMENT DEFINITION
%----------------------------------------------------------------------------------------

% article class because we want to fully customize the page and not use a cv template
\documentclass[a4paper,12pt]{article}

%----------------------------------------------------------------------------------------
%	FONT
%----------------------------------------------------------------------------------------

% % fontspec allows you to use TTF/OTF fonts directly
% \usepackage{fontspec}
% \defaultfontfeatures{Ligatures=TeX}

% % modified for ShareLaTeX use
% \setmainfont[
% SmallCapsFont = Fontin-SmallCaps.otf,
% BoldFont = Fontin-Bold.otf,
% ItalicFont = Fontin-Italic.otf
% ]
% {Fontin.otf}

%----------------------------------------------------------------------------------------
%	PACKAGES
%----------------------------------------------------------------------------------------
\usepackage{url}
\usepackage{parskip} 	

%other packages for formatting
\RequirePackage{color}
\RequirePackage{graphicx}
\usepackage[usenames,dvipsnames]{xcolor}
\usepackage[scale=0.9]{geometry}

%tabularx environment
\usepackage{tabularx}

%for lists within experience section
\usepackage{enumitem}

% centered version of 'X' col. type
\newcolumntype{C}{>{\centering\arraybackslash}X} 

%to prevent spillover of tabular into next pages
\usepackage{supertabular}
\usepackage{tabularx}
\newlength{\fullcollw}
\setlength{\fullcollw}{0.47\textwidth}

%custom \section
\usepackage{titlesec}				
\usepackage{multicol}
\usepackage{multirow}

%CV Sections inspired by: 
%http://stefano.italians.nl/archives/26
\titleformat{\section}{\Large\scshape\raggedright}{}{0em}{}[\titlerule]
\titlespacing{\section}{0pt}{10pt}{10pt}

%for publications
\usepackage[style=authoryear,sorting=ynt, maxbibnames=2]{biblatex}

%Setup hyperref package, and colours for links
\usepackage[unicode, draft=false]{hyperref}
\definecolor{linkcolour}{rgb}{0,0.2,0.6}
\hypersetup{colorlinks,breaklinks,urlcolor=linkcolour,linkcolor=linkcolour}
\addbibresource{citations.bib}
\setlength\bibitemsep{1em}

%for social icons
\usepackage{fontawesome5}

%debug page outer frames
%\usepackage{showframe}


% job listing environments
\newenvironment{jobshort}[2]
    {
    \begin{tabularx}{\linewidth}{@{}l X r@{}}
    \textbf{#1} & \hfill &  #2 \\[3.75pt]
    \end{tabularx}
    }
    {
    }

\newenvironment{joblong}[2]
    {
    \begin{tabularx}{\linewidth}{@{}l X r@{}}
    \textbf{#1} & \hfill &  #2 \\[3.75pt]
    \end{tabularx}
    \begin{minipage}[t]{\linewidth}
    \begin{itemize}[nosep,after=\strut, leftmargin=1em, itemsep=3pt,label=--]
    }
    {
    \end{itemize}
    \end{minipage}    
    }



%----------------------------------------------------------------------------------------
%	BEGIN DOCUMENT
%----------------------------------------------------------------------------------------
\begin{document}

% non-numbered pages
\pagestyle{empty} 

%----------------------------------------------------------------------------------------
%	TITLE
%----------------------------------------------------------------------------------------

% \begin{tabularx}{\linewidth}{ @{}X X@{} }
% \huge{Your Name}\vspace{2pt} & \hfill \emoji{incoming-envelope} email@email.com \\
% \raisebox{-0.05\height}\faGithub\ username \ | \
% \raisebox{-0.00\height}\faLinkedin\ username \ | \ \raisebox{-0.05\height}\faGlobe \ mysite.com  & \hfill \emoji{calling} number
% \end{tabularx}

\begin{tabularx}{\linewidth}{@{} C @{}}
\Huge{Jovanny Jara M.} \\[7.5pt]
\href{https://github.com/jovajara}{\raisebox{-0.05\height}\faGithub\ jovajara} \ $|$ \ 
\href{mailto:jvnny.jara@protonmail.com}{\raisebox{-0.05\height}\faEnvelope \ jvnny.jara@protonmail.com} \ $|$ \ {\raisebox{-0.05\height} \faMapMarked \ Av. Gran Bretaña 1111, Valparaíso, Chile} \
 \\
\end{tabularx}

%----------------------------------------------------------------------------------------
% EXPERIENCE SECTIONS
%----------------------------------------------------------------------------------------

%Interests/ Keywords/ Summary
\section{Summary}
Undergraduate Physics student (B.S. in Physics with a Minor in Scientific Computing) with demonstrated experience in Numerical Methods (C, Python) applied to complex physical systems, including General Relativity. Primary scientific interests include High Energy Physics (HEP), Cosmology, and Gravitation.
  
%Projects
\section{Projects}

\begin{tabularx}{\linewidth}{ @{}l r@{} }
July 2025 & \textbf{GeodesicSim - A Geodesic Calculator And Viewer} \hfill \href{https://github.com/jovajara/GeodesicSim}{GitHub Repository} \\[3.75pt]
\multicolumn{2}{@{}X@{}}{Final project for an elective Numerical Relativity course at Universidad de Valparaíso. Developed a C-based simulation and Python visualization tool to numerically solve time-like geodesics in Schwarzschild spacetime using the Leapfrog method.} \vspace{4mm}\\

Present & \textbf{helium-orbitals - Helium Probability Density Viewer} \hfill (In Progress) \\[3.75pt]
\multicolumn{2}{@{}X@{}}{Mandatory project for GPU Programming course (Universidad de Valparaíso). Its goal is to develop a CUDA/C-based simulation and Python visualization tool to numerically model the Helium atom orbitals using the Hartree-Fock (SCF) method.} \\
\end{tabularx}

\section{Teaching Experience}
\begin{tabularx}{\linewidth}{@{}l X@{}}
Apr 2025 - Present & \textbf{Private Tutor} (High School Physics and Mathematics) \\ 
\end{tabularx}

%----------------------------------------------------------------------------------------
%	EDUCATION
%----------------------------------------------------------------------------------------
\section{Education}
\begin{tabularx}{\linewidth}{@{}l X@{}}	
2020 - Present & Bachelor of Science (B.S.) in Physics at \textbf{Universidad de Valparaíso}, Valparaíso, Chile. \hfill \normalsize (Exp. Grad: Dec. 2026) \\
\end{tabularx}
%----------------------------------------------------------------------------------------
%	PUBLICATIONS
%----------------------------------------------------------------------------------------
% \section{Publications}
% \begin{refsection}[citations.bib]
% \nocite{*}
% \printbibliography[heading=none]
% \end{refsection}

\section{Academic Activities}
\begin{tabularx}{\linewidth}{@{}l X@{}}
Aug. 2025 - Present & Journal Club - High Energy Physics. \\ 

Jul. 2022 - Dec. 2023 & Astronomical Observation Sessions - Telescope operator and monitor.\\
\end{tabularx}

\section{Outreach}
\begin{tabularx}{\linewidth}{@{}l X@{}}

May 2023 & Verano de estrellas v2.0, Peñablanca, Valparaíso. - Telescope monitor
\end{tabularx}

%----------------------------------------------------------------------------------------
%	SKILLS
%----------------------------------------------------------------------------------------
\section{Skills}
\begin{tabularx}{\linewidth}{@{}l X@{}}
Programming Languages &  \normalsize{C (Intermediate), Python (Intermediate), CUDA (Novice), SQL (Novice), Bash (Novice).}\\
Languages  &  \normalsize{Spanish (Native), English.}\\  
\end{tabularx}

\vfill
\center{\footnotesize Last updated: \today}

\end{document}
